\subsectionaddtolist{الف}


\lr{$R^2$} یا ضریب تعیین، معیاری برای سنجش کیفیت برازش مدل رگرسیون است و نشان می‌دهد که چه نسبتی از واریانس کل خروجی‌ها توسط مدل توضیح داده شده است.

فرمول محاسبه آن به صورت زیر است:
\[
R^2 = 1 - \frac{\text{SSE}}{\text{SST}}
\]

که در آن:
\begin{itemize}
	\item $\text{SSE} = \sum_i (y^i - \hat{y}^i)^2$ \quad مجموع مربعات باقی‌مانده یا خطاها (\lr{Sum of Squared Errors})
	\item $\text{SST} = \sum_i (y^i - \bar{y})^2$ \quad مجموع مربعات کل (\lr{Total Sum of Squares})
\end{itemize}

در نتیجه، $R^2$ به ما می‌گوید که مدل چه سهمی از کل تغییرات داده‌ها را توانسته با پیش‌بینی‌هایش پوشش دهد.

\subsectionaddtolist{ب}


\begin{itemize}
	\item \textbf{مقدار \lr{$R^2$} به ۱ نزدیک می‌شود} زمانی که مدل رگرسیون بتواند تغییرات خروجی‌های واقعی را به‌خوبی پیش‌بینی کند. در این حالت، مقدار $\text{SSE}$ بسیار کوچک می‌شود و تقریباً تمام واریانس موجود در داده‌ها توسط مدل توضیح داده می‌شود:
	\[
	R^2 \approx 1 \quad \Leftrightarrow \quad \text{SSE} \ll \text{SST}
	\]
	
	\item \textbf{مقدار \lr{$R^2$} ممکن است منفی شود} اگر مدل بسیار ضعیف عمل کند و پیش‌بینی‌های آن از پیش‌بینی ساده میانگین $\bar{y}$ نیز بدتر باشد. این یعنی:
	\[
	\text{SSE} > \text{SST} \quad \Rightarrow \quad R^2 < 0
	\]
	این حالت معمولاً زمانی رخ می‌دهد که مدل اشتباه انتخاب شده باشد یا یادگیری مدل به‌درستی انجام نشده باشد.
\end{itemize}

\pagebreak

\subsectionaddtolist{ج}


\paragraph{مقایسه مفهومی:}
\begin{itemize}
	\item \lr{MSE} (میانگین مربعات خطا) نشان‌دهنده مقدار میانگین خطای مدل در واحد اندازه‌گیری خروجی است. مقدار آن همیشه مثبت است و هرچه کمتر باشد، بهتر.
	\item \lr{$R^2$} مقیاسی نسبی است که عملکرد مدل را نسبت به یک مدل ساده (که همیشه میانگین $\bar{y}$ را پیش‌بینی می‌کند) مقایسه می‌کند. این مقدار می‌تواند مثبت یا منفی باشد.
\end{itemize}

\paragraph{وضعیت تضاد:}
$\\ \\$
ممکن است یک مدل مقدار \lr{MSE} پایینی داشته باشد، ولی مقدار \lr{$R^2$} آن پایین یا حتی منفی باشد، در صورتی که:
\begin{itemize}
	\item مقادیر واقعی $y^i$ تغییرات بسیار کمی داشته باشند (یعنی $\text{SST}$ بسیار کوچک باشد).
	\item در این حالت حتی خطای کوچک مدل (مقدار $\text{SSE}$ نسبتاً کوچک) نیز ممکن است از $\text{SST}$ بزرگ‌تر باشد و $R^2$ را منفی کند.
\end{itemize}

\paragraph{مثال معکوس:}
$\\ \\$
همچنین ممکن است مقدار \lr{$R^2$} نسبتاً بالا باشد ولی مقدار \lr{MSE} بزرگ باشد؛ این زمانی رخ می‌دهد که:
\begin{itemize}
	\item داده‌های واقعی $y^i$ دامنه تغییرات بزرگی داشته باشند (یعنی $\text{SST}$ بسیار بزرگ باشد)، در نتیجه حتی اگر $\text{SSE}$ نیز بزرگ باشد، نسبت $\frac{\text{SSE}}{\text{SST}}$ کوچک باقی بماند و $R^2$ عدد بزرگی شود.
\end{itemize}



\subsectionaddtolist{د}


برای به‌دست آوردن گرادیان، باید مشتق جزئی تابع خطا را نسبت به هر پارامتر بگیریم.

{گرادیان نسبت به \boldmath$w_2$:}
\[
\frac{\partial f}{\partial w_2} = \sum_i 2 \left( \hat{y}^i - y^i \right) \cdot (x^i)^2
\]

{گرادیان نسبت به \boldmath$w_1$:}
\[
\frac{\partial f}{\partial w_1} = \sum_i 2 \left( \hat{y}^i - y^i \right) \cdot x^i
\]

{گرادیان نسبت به \boldmath$w_0$:}
\[
\frac{\partial f}{\partial w_0} = \sum_i 2 \left( \hat{y}^i - y^i \right)
\]

\pagebreak


\subsectionaddtolist{ه}


در روش \lr{Gradient Descent}، پارامترهای مدل به‌صورت تکراری به سمت خلاف جهت گرادیان تابع خطا به‌روزرسانی می‌شوند تا مقدار تابع خطا کاهش یابد. رابطه کلی به شکل زیر است:

\[
\theta \leftarrow \theta - \alpha \cdot \nabla_\theta f(\theta)
\]

که در آن:
\begin{itemize}
	\item $\theta$ نمایان‌گر پارامترهای مدل (مثلاً $w_2, w_1, w$) است.
	\item $\nabla_\theta f(\theta)$ گرادیان تابع خطا نسبت به $\theta$ است.
	\item $\alpha$ نرخ یادگیری است.
\end{itemize}

\paragraph{تأثیر تغییر $\alpha$:}

\begin{itemize}
	\item اگر $\alpha$ خیلی کوچک باشد:  
	به‌روزرسانی پارامترها کند انجام می‌شود و همگرایی مدل بسیار آهسته خواهد بود.
	
	\item اگر $\alpha$ خیلی بزرگ باشد:  
	به‌روزرسانی‌ها بیش‌از‌حد شدید شده و ممکن است مدل از نقطه بهینه عبور کند یا دچار نوسان و واگرایی شود.
	
	\item انتخاب مناسب $\alpha$:  
	نرخ یادگیری باید به اندازه‌ای باشد که مدل سریع و پایدار به سمت کمینه تابع خطا حرکت کند.
\end{itemize}
